%% ====================================================================
\section{Results}
\label{sec:results}

\subsection{Does the Sensor-Level Bound Hold? Same-Window Calibration}
\label{sec:bound}

This subsection is the setting of Theorem~\ref{thm:coverage}: the calibration sensors' errors are those of the window being evaluated. It is a reference for what the theorem delivers, not a deployable procedure; Section~\ref{sec:deployable} removes that advantage.

Table~\ref{tab:bound} compares the two calibration rules at the 90\% target. For each trained model
we draw 4000 random (new sensor, calibration set) assignments from $\mathcal{R}$ (1500 for the two
traffic networks) and average the new sensor's coverage; the table gives the mean over splits.
Where Theorem~\ref{thm:coverage} gives a finite interval, its average coverage is at or above 0.90 in
every case, from 0.911 (PeMS-BAY, $n{=}89$) to 0.9996 (Delhi, $n{=}9$). The same holds at the 80\%
target (0.810 to 0.888 with at least 9 calibration sensors; the full set is in the released analysis
file). This is what the theorem predicts, and it is conservative: the guarantee costs width.

\begin{table}[htbp]
\caption{Coverage and Width at the 90\% Target, Sensor-Level Rule Versus Pooled Rule}
\label{tab:bound}
\centering
\footnotesize
\setlength{\tabcolsep}{3pt}
\begin{tabular}{lrcccccc}
\toprule
 & & \multicolumn{2}{c}{\textbf{Sensor-level}} & \multicolumn{2}{c}{\textbf{Pooled}} & \textbf{Width} & \textbf{Draws} \\
\textbf{Dataset} & $n$ & Cov. & Width & Cov. & Width & \textbf{ratio} & \textbf{$<0.80$} \\
\midrule
Delhi AQ   &  9 & 0.9996 & 491 & 0.8998 & 127.6 & 3.85 & --- \\
           & 15 & 0.9597 & 204 & 0.8997 & 127.6 & 1.60 & 0\% \\
Beijing AQ &  5 & vacuous & $\infty$ & 0.9003 & 179.0 & --- & 0\% \\
Weather    &  9 & 0.9996 & 36.8 & 0.8960 & 20.3 & 1.82 & --- \\
           & 11 & 0.9795 & 27.8 & 0.8965 & 20.3 & 1.37 & 13\% \\
PeMSD7-M   & 20 & 0.9443 & 27.3 & 0.8987 & 17.6 & 1.55 & --- \\
           & 59 & 0.9155 & 20.1 & 0.9003 & 17.6 & 1.14 & 5\% \\
PeMS-BAY   & 20 & 0.9448 & 40.8 & 0.8983 & 19.4 & 2.10 & --- \\
           & 89 & 0.9110 & 21.3 & 0.9010 & 19.4 & 1.10 & 8\% \\
\bottomrule
\multicolumn{8}{p{0.93\columnwidth}}{\footnotesize Means over 10 splits (Delhi, Beijing, weather) or 4
(traffic); widths in the units of each dataset. $n$ is the number of calibration sensors. ``Draws
$<0.80$'' is the share of random assignments, under the pooled rule at the largest $n$, whose new
sensor was covered less than 0.80 of the time. With 6 held-out sensors (Beijing) the sensor-level rule
is vacuous at the 90\% target; at the 80\% target and $n{=}5$ it gives coverage 0.960 at width 310
against 0.801 at width 129 for the pooled rule.} \\
\end{tabular}
\end{table}

\begin{figure*}[!t]
\centering
\includegraphics[width=0.72\textwidth]{fig_sensor_level.png}
\caption{Mean coverage as the number of calibration sensors grows, at the 90\% target (dotted line),
same-window calibration. The sensor-level rule stays above target and its price (Table~\ref{tab:bound}) falls as
$n$ grows; the pooled rule sits on the target on average.}
\label{fig:sensorlevel}
\end{figure*}

Three points follow. \emph{First}, the price is set by the number of calibration sensors and not by
the number of calibration points: with 9 to 15 sensors the interval is 1.4 to 4.3 times wider than the
pooled one, with 59 to 89 sensors it is 1.1 times wider. A network with a dozen held-out sensors cannot
support a finite distribution-free 90\% interval at all. \emph{Second}, the pooled rule, which has no
guarantee here, averages between 0.896 and 0.901 over random new sensors at every $n$ we tried. We read
this as an empirical fact about these networks, not as a property we can promise: in the weather data
it is 0.4 points under target, and the average hides the spread. For the same pooled rule, in 13\%
(weather), 8\% (PeMS-BAY) and 5\% (PeMSD7-M) of random assignments the new sensor is covered less
than 0.80 of the time, and the tenth percentile of the per-draw coverage is 0.78 to 0.88. Marginal
coverage is a statement about the average new sensor. \emph{Third}, nothing in Theorem~\ref{thm:coverage}
depends on OmniST: it applies to any forecaster whose error at a new sensor is a function of that
sensor's data and the training network, including the baselines below.

\subsection{Calibration at Forecast Time}
\label{sec:deployable}

Table~\ref{tab:deployable} repeats the evaluation with the threshold for each window formed only from
errors of earlier windows (Section~\ref{sec:online-conformal}), at the largest calibration size of each
network and the 90\% target. The same random (new sensor, calibration set) draws are used for every scheme.

\begin{table*}[!t]
\caption{Coverage at Forecast Time (Target 0.90). Static, Expanding, Rolling and ACI Calibration, Sensor-Level and Pooled Rules, Largest $n$}
\label{tab:deployable}
\centering
\footnotesize
\setlength{\tabcolsep}{5pt}
\begin{tabular}{lrlcccccc}
\toprule
 & & & \multicolumn{4}{c}{\textbf{Mean coverage}} & \multicolumn{2}{c}{\textbf{Share of draws $<0.80$}} \\
\textbf{Dataset} & $n$ & \textbf{Rule} & Static & Expanding & Rolling & ACI & Rolling & ACI \\
\midrule
Delhi AQ   & 15 & Sensor-level & 0.848 & 0.863 & 0.853 & 0.855 & 0.10 & 0.02 \\
           &    & Pooled       & 0.800 & 0.795 & 0.785 & 0.790 & 0.55 & 0.53 \\
Weather    & 11 & Sensor-level & 0.948 & 0.970 & 0.971 & 0.947 & 0.02 & 0.00 \\
           &    & Pooled       & 0.854 & 0.882 & 0.889 & 0.863 & 0.13 & 0.19 \\
PeMSD7-M   & 59 & Sensor-level & 0.914 & 0.914 & 0.913 & 0.900 & 0.05 & 0.00 \\
           &    & Pooled       & 0.900 & 0.899 & 0.897 & 0.883 & 0.07 & 0.00 \\
PeMS-BAY   & 89 & Sensor-level & 0.870 & 0.890 & 0.899 & 0.899 & 0.04 & 0.00 \\
           &    & Pooled       & 0.854 & 0.877 & 0.889 & 0.887 & 0.12 & 0.00 \\
\bottomrule
\multicolumn{9}{p{0.9\textwidth}}{\footnotesize Means over splits (10 for Delhi and weather, 4 for traffic) of
the new sensor's coverage averaged over the evaluated windows. Evaluation starts at window $L$: 6 windows for
Delhi, 18 for weather, 183 for PeMSD7-M and 116 for PeMS-BAY. Rolling uses $L=4$ windows (daily networks) or 24
(traffic); ACI uses $\gamma=0.1$. No theorem covers these schemes.} \\
\end{tabular}
\end{table*}

Four statements follow. \emph{First}, the same-window result is an upper bound on what the theorem-backed
procedure delivers at forecast time. With static calibration the sensor-level rule falls to 0.870 on
PeMS-BAY ($n{=}89$) and 0.848 on Delhi. \emph{Second}, rolling recalibration recovers the target on PeMS-BAY
(0.899) and keeps weather (0.971) and PeMSD7-M (0.913) above it, but not on Delhi, where no scheme exceeds
0.863. Delhi has six evaluation windows of 30 days and a strong seasonal cycle, and we do not claim the other
networks are free of drift. \emph{Third}, the pooled rule, which was on target in the same-window setting,
under-covers on Delhi, weather and PeMS-BAY under every scheme (0.785 to 0.889), and on PeMSD7-M it is within
0.003 of the target except under ACI (0.883). \emph{Fourth}, ACI removes the tail of badly covered sensors: the
share of random new sensors covered below 0.80 under the sensor-level rule falls from 10\%, 2\%, 5\% and 4\%
(rolling) to 2\%, 0\%, 0\% and 0\%. It does so by changing the width: relative to rolling, ACI intervals are 7\%
narrower on Delhi, 15\% on weather and 3\% on PeMSD7-M, and 31\% wider on PeMS-BAY (28.8 against 21.9). The
earlier version of this work reported that a rolling ACI layer raised late-year coverage on Delhi from 0.50 to
0.96. That result came from a different protocol and does not reproduce here; we withdraw it.

\subsection{What the Earlier Protocol Would Have Reported}
\label{sec:legacy}

The previous version of this work calibrated on the residuals of the training sensors and measured
coverage on 3 to 5 fixed held-out sensors, the same in every seed. It reported 0.914 to 0.928.
Table~\ref{tab:legacy} repeats the evaluation with sensors split at random into training,
calibration and test sets, over many splits (experiment C1, released with the code), once with disjoint
calibration sensors and once, on the same trained models, with training-sensor calibration. Mean
coverage is 0.876 to 0.906; three of five estimates lie below 0.90, and the earlier numbers were 1.3 to
3.9 points higher than the random-split means. On the same trained models, calibrating on training
sensors and calibrating on disjoint sensors differ by at most 0.018 in average coverage, in either
direction. The concern about exchangeability therefore did not turn into a large average error on
these data, but it remained unsupported until the sensor-level argument above.

\begin{table}[htbp]
\caption{Coverage Over Random Disjoint Sensor Splits (Target 0.90, Pooled Rule)}
\label{tab:legacy}
\centering
\footnotesize
\setlength{\tabcolsep}{3pt}
\begin{tabular}{lrccc}
\toprule
 & & \textbf{Disjoint cal.} & \textbf{Train-sensor} & \textbf{Worst} \\
\textbf{Dataset} & \textbf{Splits} & \textbf{(95\% interval)} & \textbf{cal.} & \textbf{sensor} \\
\midrule
Delhi AQ   & 20 & 0.893 [0.884, 0.901] & 0.901 & 0.777 \\
Beijing AQ & 20 & 0.903 [0.895, 0.912] & 0.902 & 0.861 \\
Weather    & 20 & 0.906 [0.881, 0.931] & 0.888 & 0.654 \\
PeMSD7-M   &  6 & 0.876 [0.843, 0.909] & 0.890 & 0.519 \\
PeMS-BAY   &  6 & 0.892 [0.872, 0.913] & 0.887 & 0.402 \\
\bottomrule
\multicolumn{5}{p{0.93\columnwidth}}{\footnotesize ``Worst sensor'' is the lowest per-sensor
coverage over all test sensors and splits, under disjoint calibration.} \\
\end{tabular}
\end{table}

\subsection{Accuracy at New Sensors and What Each Component Contributes}
\label{sec:ablation}

\begin{table}[htbp]
\caption{MAE at Held-Out Sensors, Mean Over Splits (95\% Interval)}
\label{tab:accuracy}
\centering
\footnotesize
\setlength{\tabcolsep}{3.2pt}
\begin{tabular}{lccccc}
\toprule
\textbf{Method} & \textbf{Delhi} & \textbf{Beijing} & \textbf{Weather} & \textbf{D7-M} & \textbf{BAY} \\
\midrule
OmniST (normalized)    & 28.5 & 48.8 & 4.75 & 3.66 & 4.49 \\
No attention           & 28.6 & 48.8 & 4.76 & ---  & ---  \\
Coordinate-free GRU    & 28.0 & 48.4 & 4.82 & 3.71 & 4.45 \\
Persistence            & 38.0 & 52.7 & 4.65 & 3.57 & 5.14 \\
Window mean            & 47.2 & 51.3 & 4.58 & 7.44 & 5.41 \\
OmniST, earlier code   & 28.0 & 49.6 & 5.72 & 9.76 & 5.17 \\
\midrule
IDW, concurrent$^\dagger$ & 18.9 & 8.4 & 3.22 & 7.63 & 5.66 \\
\bottomrule
\multicolumn{6}{p{0.93\columnwidth}}{\footnotesize Units of each dataset; D7-M is PeMSD7-M, BAY is
PeMS-BAY. ``Normalized'': coordinates normalized from the training network. ``Earlier code'' uses the
coordinate constants of the submitted version (see text). $^\dagger$Inverse-distance weighting of the
training sensors' \emph{concurrent} values: a nowcasting reference that a forecaster cannot use.} \\
\end{tabular}
\end{table}

Table~\ref{tab:accuracy} and the paired differences (95\% intervals over splits) lead to four
statements.

\emph{(i) Against no-training forecasters.} OmniST is clearly better than persistence on the
air-quality networks (Delhi $-9.5$ [$-10.8$, $-8.1$]; Beijing $-3.9$ [$-4.3$, $-3.5$]) and on PeMS-BAY
($-0.65$ [$-0.76$, $-0.55$]), and indistinguishable from it on weather ($+0.10$ [$-0.23$, $0.43$]) and
PeMSD7-M ($+0.09$ [$-0.19$, $0.37$]). On those two networks a forecaster that repeats the last value
is as accurate. The comparison with the concurrent-observation interpolators is not one of equals:
they see what happens at the neighboring sensors during the target window and they are better
on Delhi, Beijing and weather. We therefore withdraw the earlier statement that OmniST improves on
interpolation by 12\% and on a kriging network by 21\%, which compared a forecast with a
reconstruction.

\emph{(ii) The coordinate-free control.} A GRU shared across sensors, with neither coordinates nor
attention, is within about 2\% of OmniST on every dataset: OmniST minus control is $+0.57$ [$-0.20$, $1.33$]
on Delhi, $+0.42$ [$-0.02$, $0.87$] on Beijing, $-0.067$ [$-0.122$, $-0.012$] on weather, $-0.056$
[$-0.43$, $0.32$] on PeMSD7-M and $+0.04$ [$-0.03$, $0.11$] on PeMS-BAY. Removing only the attention
changes the error by at most 0.11 on the datasets where we ran it. \textbf{The capability to forecast
at a new sensor comes from sharing parameters across sensors, not from the coordinates or the
attention.} In the data we have, the coordinates and the attention do not improve accuracy at new
sensors; they make spatial context available, and the architecture keeps the locality property
that Lemma~\ref{lem:locality} needs, but we cannot claim a benefit from them.

\emph{(iii) A defect in the earlier code.} The coordinate embedding in the submitted version centered
and scaled coordinates with constants chosen for Delhi (latitude 28.6, longitude 77.1, scales 0.2 and
0.3). For any other region the input to the embedding MLP lay tens to hundreds of units outside its
intended range, and the control experiment exposed it: with those constants OmniST was 0.7 to 6.0
error units worse than the coordinate-free control on every dataset but Delhi (PeMSD7-M: 9.76 against
3.71). Normalizing coordinates with the mean and standard deviation of the \emph{training} network,
fixed before any held-out sensor is seen, removes the gap (PeMSD7-M: 3.66; weather: 5.72 to 4.75;
PeMS-BAY: 5.17 to 4.49). Delhi, whose range the constants were chosen for, is slightly worse after the
change ($+0.48$ [$0.24$, $0.73$]). Every coverage and accuracy result in this paper uses the
normalized version; Table~\ref{tab:accuracy} keeps the earlier values for the record.

\emph{(iv) Intervals are as wide as the forecaster is inaccurate.} The width of every interval follows
the error of its forecaster: persistence gives intervals about 1.5 times wider than OmniST on Delhi
and PeMS-BAY, and nearest-neighbor interpolation gives intervals about 2.9 times wider on PeMSD7-M, at
the same coverage.
Calibration, not the forecaster, is what makes the coverage hold.

\emph{(v) Against seasonal-naive.} A seasonal-naive forecast repeats the sensor's own value one period earlier
(a week for the daily networks, the previous day for the traffic networks), restricted for all methods to the
windows that have one. It is no better than persistence on Delhi and clearly worse than OmniST on Delhi
($-9.5$ [$-10.7$, $-8.3$], 10 of 10 splits) and Beijing ($-15.2$ [$-15.9$, $-14.5$]); it ties OmniST on weather
($-0.24$ [$-0.56$, $0.07$]) and is worse on PeMSD7-M ($-0.74$ [$-1.16$, $-0.32$]). On PeMS-BAY it is clearly better
than both learned models: MAE 3.56 against 4.68 for OmniST (OmniST minus seasonal-naive $+1.12$ [$0.95$, $1.30$],
0 of 4 splits), and its 90\% pooled interval is 15.6 wide against 20.5, so OmniST's intervals are 31\% wider.
\emph{(vi) Removing reference sensors} (30\% or 50\% of the training network at inference, rescoring and
recalibrating on the survivors) changes MAE by at most 0.6\% and coverage by about 0.3 points at most
(supplementary material). This is expected: the model reads only its $k$ nearest neighbors and, as the controls
show, gets little from them.

\subsection{Distance From the Training Network}
\label{sec:distance}

Assumption (A1) is least safe for a sensor far from the network. Holding out a contiguous region and relating
per-sensor coverage to the distance from the nearest training sensor, in units of the median spacing
(experiment C2), a distance effect appears only in the sparse weather network: coverage falls by 0.0136 per extra
spacing [$-0.0221$, $-0.0063$]. In the other four networks the within-sensor slopes lie within 0.004 of zero. A
total-variation bound on the shift between sensors' score distributions would be vacuous at our sample sizes
(experiment C3: median 0.11 on Delhi and 0.19 on weather, against a noise floor of 0.13). Weather is also where the
pooled rule most often misses badly for an individual sensor, but with one sparse network of 30 sensors this is a
pattern to test, not a result.

@@LIMITS@@
